\chapter{Gładkie przeszkody Donaldsona w wymiarze cztery}
\label{ch:donaldson}
\index{twierdzenie Donaldsona}
\index{równanie antysamodualności}

Klasyfikacja Freedmana z Rozdziału~\ref{ch:freedman-klasyfikacja}
dotyczy kategorii topologicznej. Pytanie, które z jej przykładów
można wygładzić, wymaga innego narzędzia. Donaldson wykorzystał
połączenia na wiązkach nad gładką rozmaitością czterowymiarową
i analizę ich krzywizny. Poniżej wyprowadzamy równanie
antysamodualności oraz jego prosty rachunek energii. Głębokie
twierdzenie o przestrzeni rozwiązań i wynikającą z niego
diagonalizację podajemy jako wynik zewnętrzny.

\section{Krzywizna połączenia i gwiazda Hodge'a}
\label{sec:donaldson-asd}

Niech $X$ będzie zamkniętą, zorientowaną, gładką
rozmaitością Riemanna wymiaru $4$, a $P\to X$ główną
wiązką o grupie $\mathrm{SU}(2)$. Grupę i jej algebrę
opisuje Dodatek~\ref{sec:app-grupa-u}; połączenia
omawialiśmy w Rozdziale~\ref{ch:koneksje}, a gwiazdę Hodge'a
w Sekcji~\ref{sec:hodge-four-manifolds}.
W lokalnej trywializacji połączenie opisuje
$1$-forma macierzowa $A$ o wartościach w
$\mathfrak{su}(2)$, a jego krzywiznę definiuje wzór
\begin{equation}\label{eq:donaldson-curvature}
 F_A=dA+A\wedge A.
\end{equation}
Iloczyn w drugim składniku łączy iloczyn zewnętrzny form
z mnożeniem macierzy. Dla pól stycznych $u,v$ mamy więc
$(A\wedge A)(u,v)=[A(u),A(v)]$.
Zmiana lokalnej trywializacji nad otwartym $U\subset X$
przez $g:U\to\mathrm{SU}(2)$ daje
$A^g=g^{-1}Ag+g^{-1}dg$ oraz
$F_{A^g}=g^{-1}F_Ag$. Ostatnia równość wynika po
rozwinięciu $d(g^{-1})=-g^{-1}(dg)g^{-1}$:
wyrazy z $dg$ znoszą się parami. Dlatego warunek
na krzywiznę nie zależy od wybranej lokalnej
trywializacji także wtedy, gdy $P$ nie jest trywialna.

W wymiarze cztery gwiazda Hodge'a spełnia
$*^2=1$ na $2$-formach. Dzielimy więc krzywiznę na
\[
 F_A^+=\tfrac12(F_A+*F_A),\qquad
 F_A^-=\tfrac12(F_A-*F_A).
\]

\begin{definicja}[Połączenie antysamodualne]
\label{def:donaldson-asd}
Połączenie $A$ na $P$ jest \emph{antysamodualne},
gdy $F_A^+=0$, czyli $*F_A=-F_A$.
Przestrzeń takich połączeń modulo działanie
gładkich przekształceń cechowania nazywamy
\emph{przestrzenią moduli ASD}.
\end{definicja}

Intuicja jest następująca: gwiazda rozdziela
krzywiznę na dwie niezależne połowy, a równanie
zeruje jedną z nich. Na wiązce trywialnej połączenie
$A=0$ jest płaskie, więc stanowi prosty przykład
rozwiązania.
Nie twierdzimy, że na dowolnej wiązce istnieje
rozwiązanie niepłaskie; zależy to od topologii
wiązki i metryki.

\begin{przyklad}[Niepłaskie rozwiązanie ASD na $\R^4$]
\label{prz:donaldson-asd-lokalny}
Na standardowo zorientowanej euklidesowej $\R^4$
weźmy trywialną wiązkę o grupie $\mathrm{SU}(2)$,
$T=\operatorname{diag}(i,-i)\in\mathfrak{su}(2)$
oraz połączenie opisane globalnie przez
\[
 A=T(x^1\dd x^2-x^3\dd x^4).
\]
Ponieważ wszystkie wartości $A$ są wielokrotnościami
tej samej macierzy, $A\wedge A=0$. Zatem
\[
 F_A=\dd A
   =T(\dd x^1\wedge\dd x^2
       -\dd x^3\wedge\dd x^4).
\]
Dla orientacji $\dd x^1\wedge\cdots\wedge\dd x^4$
gwiazda zamienia $\dd x^1\wedge\dd x^2$
z $\dd x^3\wedge\dd x^4$. Stąd $*F_A=-F_A$,
choć $F_A\ne0$: jest to rozwiązanie ASD, które
nie jest płaskie. Ponieważ $|T|^2=2$ i obie
skalarne $2$-formy są ortonormalne, $|F_A|^2=4$
w każdym punkcie. Jego energia na całej $\R^4$
jest więc nieskończona; poniższa tożsamość
całkowa dotyczy zamkniętej rozmaitości $X$.
\end{przyklad}

\section{Energia i liczba Cherna}
\label{sec:donaldson-energia}

Na $\mathfrak{su}(2)$ przyjmujemy dodatni iloczyn
$\langle B,C\rangle=-\operatorname{tr}(BC)$.
Niech $\|F_A\|_{L^2}^2=\int_X|F_A|^2\,d\mathrm{vol}$.
Niech $E=P\times_{\mathrm{SU}(2)}\C^2$ będzie
stowarzyszoną wiązką zespoloną. Przy konwencji,
w której formą reprezentującą łączną klasę Cherna
jest $c_A(E)=\det(I+\frac{i}{2\pi}F_A)$.
Ponieważ $\mathrm{SU}(2)$ działa z wyznacznikiem $1$,
wiązka $\det E$ jest trywialna i $c_1(E)=0$ także
jako klasa całkowita. Równość $\operatorname{tr}F_A=0$
odzwierciedla to na poziomie form Cherna--Weila.
Współczynnik stopnia cztery w rozwinięciu wyznacznika
ma znak dodatni:
\[
 \frac12\left(\frac{i}{2\pi}\right)^2
 \bigl((\operatorname{tr}F_A)^2
       -\operatorname{tr}(F_A\wedge F_A)\bigr)
 =\frac{1}{8\pi^2}\operatorname{tr}(F_A\wedge F_A).
\]
Dlatego druga liczba Cherna wynosi
\begin{equation}\label{eq:donaldson-c2}
 k=c_2(E)[X]
   =\frac{1}{8\pi^2}\int_X\operatorname{tr}(F_A\wedge F_A)
   \in\Z.
\end{equation}
Całkowitość i niezależność od połączenia są tu
zastosowaniem teorii Cherna--Weila z
Rozdziału~\ref{ch:chern-weil} oraz identyfikacji
otrzymanej formy z całkowitą klasą $c_2(E)$.
Tam wzór na pierwszą klasę Pontriagina wiązki
\emph{rzeczywistej} ma przeciwny znak przy śladzie;
tutaj znak wynika z czynnika $i^2=-1$ w wyznaczniku.
Sam rachunek energii możemy przeprowadzić bez
analizy przestrzeni moduli.

\begin{stwierdzenie}[Rachunek energii ASD]
\label{prop:donaldson-energy}
Dla powyższych konwencji
\[
 \|F_A\|_{L^2}^2
 =8\pi^2k+2\|F_A^+\|_{L^2}^2.
\]
W szczególności dla połączenia antysamodualnego
energia wynosi $8\pi^2k$ i $k\geq0$.
\end{stwierdzenie}
\begin{proof}
Ponieważ $F_A^+$ i $F_A^-$ są ortogonalne,
\[
 |F_A|^2=|F_A^+|^2+|F_A^-|^2.
\]
W zorientowanej bazie ortonormalnej dla skalarnej
$2$-formy $\omega$ zachodzi
$\omega\wedge\omega=\langle\omega,*\omega\rangle
\,d\mathrm{vol}$. Stosując tę równość składowo
do macierzy i znak z definicji iloczynu na
$\mathfrak{su}(2)$, otrzymujemy
\[
 \operatorname{tr}(F_A\wedge F_A)
 =\bigl(|F_A^-|^2-|F_A^+|^2\bigr)
   \,d\mathrm{vol}.
\]
Równość~\eqref{eq:donaldson-c2} daje więc
$8\pi^2k=\|F_A^-\|_{L^2}^2-\|F_A^+\|_{L^2}^2$.
Dodanie $2\|F_A^+\|_{L^2}^2$ kończy dowód.
\end{proof}

Jeżeli $k=0$, każde połączenie ASD na zamkniętym
$X$ jest płaskie: jego energia wynosi zero, więc
ciągła krzywizna znika wszędzie. Przy $k<0$
rozwiązanie ASD nie może istnieć. Dlatego
niepłaski model z
Przykładu~\ref{prz:donaldson-asd-lokalny},
którego energia na całej $\R^4$ jest nieskończona,
nie przeczy temu ograniczeniu.

\section{Twierdzenie o diagonalizacji i jego zakres}
\label{sec:donaldson-diagonalizacja}

Przestrzeń moduli ASD nie jest zwykłą rozmaitością
zadaną skończoną liczbą równań. Trzeba wybrać
warunek cechowania, uzyskać lokalny model eliptyczny,
zbadać rozwiązania redukowalne oraz utratę zwartości
przez koncentrację krzywizny. Zwartość Uhlenbeck,
usuwanie osobliwości i analiza końców przestrzeni moduli
są odrębnymi głębokimi wynikami. Rozdział nie zawiera
ich dowodów.

\begin{twierdzenie}[Donaldson: diagonalizacja formy określonej;
wynik zewnętrzny]
\label{thm:donaldson-diagonalizacja}
Niech $X$ będzie zamkniętą, zorientowaną, gładką,
po prostu spójną rozmaitością wymiaru $4$.
Jeśli całkowita forma przecięcia $Q_X$ na
$H_2(X;\Z)/\operatorname{Tors}$ jest dodatnio określona,
to jest całkowicie izometryczna z
$\langle1\rangle^{\oplus b_2(X)}$.
Jeśli jest ujemnie określona, jest izometryczna z
$\langle-1\rangle^{\oplus b_2(X)}$.
\end{twierdzenie}

Pierwszą wersję otrzymuje się z drugiej po zmianie
orientacji. Źródłem jest
\href{https://doi.org/10.4310/jdg/1214437665}
{S.~K.~Donaldson, \emph{An application of gauge theory to
four-dimensional topology}, J.~Differential Geom. 18 (1983),
twierdzenie~A}. W oryginalnej orientacji dowód bada
połączenia \emph{samodualne} na wiązce z liczbą
$c_2(E)[X]=-1$ nad rozmaitością o dodatnio określonej
formie. Po odwróceniu orientacji gwiazda zmienia znak
na $2$-formach, podobnie jak ewaluacja $c_2(E)[X]$.
Otrzymujemy więc opisaną tutaj konwencję ASD z $k=1$
nad rozmaitością o ujemnie określonej formie.

Dla zrozumienia metody załóżmy teraz, że $Q_X$ jest
ujemnie określona. Rozwiązanie redukowalne odpowiada
rozkładowi $E\cong L\oplus L^{-1}$, dla którego
$1=c_2(E)[X]=-c_1(L)^2[X]$. Takie rozwiązania
odpowiadają zatem parom klas $\{a,-a\}$
o $Q_X(a,a)=-1$. Niech $n(Q_X)$ będzie liczbą tych par.
Odwrotnie, harmoniczny reprezentant każdej takiej klasy
jest antysamodualny, ponieważ forma przecięcia jest
ujemnie określona. Daje on połączenie na $L$ o krzywiźnie
ASD, a $H^1(X;\R)=0$ zapewnia jednoznaczność modulo
cechowanie. Zamiana $L$ na $L^{-1}$ nie zmienia punktu
redukowalnego.
Głęboka część dowodu Donaldsona konstruuje
pięciowymiarową przestrzeń moduli z końcem od
koncentracji krzywizny i punktami redukowalnymi;
rachunek sygnatury jej odpowiednio obciętej wersji
daje $n(Q_X)=b_2(X)$. Ten krok, w tym zwartość,
orientowalność i zachowanie przestrzeni moduli przy
punktach redukowalnych, pozostaje wynikiem zewnętrznym;
por. także \href{https://doi.org/10.1090/S0273-0979-1983-15090-5}
{Donaldson, \emph{Self-dual connections and the topology
of smooth 4-manifolds}, \S\,2--4}.

Końcowy krok jest już algebraiczny. Dla różnych par
$\{a,-a\}$ i $\{b,-b\}$ o kwadracie $-1$ nierówność
Cauchy'ego--Schwarza dla $-Q_X$ daje
$|Q_X(a,b)|\leq1$. Równość oznaczałaby $a=\pm b$,
więc całkowitość formy wymusza $Q_X(a,b)=0$.
Przedstawiciele $n(Q_X)=b_2(X)$ par tworzą zatem
ortogonalny układ pełnej rangi o macierzy $-I$.
Wyznacznik ograniczenia $Q_X$ do generowanej przez
nie podgrupy ma wartość bezwzględną $1$, tak jak
wyznacznik formy $Q_X$. Kwadrat indeksu tej podgrupy
wynosi więc $1$, co czyni ten układ bazą całego
$H_2(X;\Z)/\operatorname{Tors}$. Nie zastępuje to
analitycznej części dowodu Donaldsona.

\begin{wniosek}[Forma $E_8$ nie pochodzi od gładkiej
zamkniętej jednospójnej $4$-rozmaitości]
\label{cor:donaldson-e8}
Nie istnieje zamknięta, gładka, po prostu spójna
rozmaitość $4$-wymiarowa o formie przecięcia $E_8$.
Topologiczna rozmaitość Freedmana o tej formie
nie dopuszcza struktury gładkiej.
\end{wniosek}
\begin{proof}
Z obliczenia w
Stwierdzeniu~\ref{prop:freedman-E8-properties} forma
$E_8$ jest dodatnio określona, parzysta i ma rangę $8$.
Twierdzenie~\ref{thm:donaldson-diagonalizacja}
wymuszałoby izometrię z $I_8$. Ta ostatnia forma
jest nieparzysta, bo wektor bazowy ma kwadrat $1$;
parzystość zachowują izometrie całkowite. To sprzeczność.
Istnienie topologicznej rozmaitości o tej formie jest
częścią zewnętrznego twierdzenia Freedmana przyjętego
w Rozdziale~\ref{ch:freedman-klasyfikacja}.
\end{proof}

\begin{wniosek}[Zniknięcie przeszkody stabilnej
nie wystarcza do wygładzenia]
\label{cor:donaldson-2e8}
Topologiczna, zamknięta, zorientowana, jednospójna
rozmaitość Freedmana $X_{2E_8}$ o formie
$E_8\oplus E_8$ spełnia $\operatorname{ks}(X_{2E_8})=0$,
lecz nie dopuszcza struktury gładkiej.
\end{wniosek}
\begin{proof}
Stwierdzenie~\ref{prop:freedman-E8-properties}
pokazuje, że $E_8$ jest parzysta, unimodularna
i dodatnio określona oraz że $\sigma(E_8)=8$.
Suma prosta dwóch kopii jest więc parzystą,
unimodularną formą dodatnio określoną rangi $16$
o sygnaturze $16$. Twierdzenie
\ref{thm:freedman-classification} daje rozmaitość
topologiczną o tej formie, a warunek
\eqref{eq:freedman-even-ks} daje
$\operatorname{ks}(X_{2E_8})\equiv16/8=2\equiv0\pmod2$,
zgodnie z Przykładem~\ref{prz:freedman-examples}.
Gdyby $X_{2E_8}$ była gładka, Twierdzenie
\ref{thm:donaldson-diagonalizacja} wymagałoby
izometrii $E_8\oplus E_8\cong I_{16}$ nad $\Z$.
Lewa forma jest parzysta, natomiast $I_{16}$
jest nieparzysta, bo wektor bazowy ma kwadrat $1$.
Izometria zachowuje parzystość, więc taka
izometria nie istnieje.
\end{proof}

To porównanie wyostrza Wniosek~\ref{cor:donaldson-e8}:
tam $\operatorname{ks}(X_{E_8})=1$ już wyklucza
gładkość, natomiast tutaj przeszkoda stabilna znika.
Inny wynik zewnętrzny, twierdzenie Rokhlina, mówi,
\index{twierdzenie Rokhlina}
że sygnatura każdej zamkniętej, zorientowanej,
gładkiej rozmaitości spinowej wymiaru cztery
jest podzielna przez $16$;
zob. \href{https://arxiv.org/pdf/2011.12418}
{M.~R.~Klug, \emph{A relative version of Rochlin's theorem},
wprowadzenie}. Gdyby $X_{E_8}$ była gładka,
jej parzysta forma i wzór Wu z
Sekcji~\ref{sec:freedman-ks} czyniłyby ją spinową,
ale sygnatura $8$ naruszałaby ten warunek.
Dla $X_{2E_8}$ sygnatura wynosi $16$, więc
twierdzenie Rokhlina nie wyklucza wygładzenia;
rozstrzyga je dopiero diagonalizacja Donaldsona.

Oba wnioski pokazują precyzyjną granicę między
klasyfikacją topologiczną a gładką. Nie wynika z nich
jeszcze istnienie egzotycznej $\R^4$: potrzebna jest
osobna konstrukcja otwartej rozmaitości i zwarta
przeszkoda rozróżniająca struktury. Przeprowadzimy
wynikającą z niej indukcję w następnym rozdziale.

\par\smallskip
{\footnotesize\noindent\emph{Dalsza lektura:}
\href{https://doi.org/10.4310/jdg/1214437665}
{S.~K.~Donaldson, \emph{An application of gauge theory to
four-dimensional topology}, J.~Differential Geom. 18 (1983)},
twierdzenie o diagonalizacji;
\href{https://doi.org/10.1090/S0273-0979-1983-15090-5}
{S.~K.~Donaldson, \emph{Self-dual connections and the topology
of smooth 4-manifolds}, Bull.~Amer.~Math.~Soc. 8 (1983)},
metoda cechowania i kontekst wyniku.\par}
